\section*{ВИСНОВКИ ТА ПЕРСПЕКТИВИ ПОДАЛЬШИХ ДОСЛІДЖЕНЬ}

У статті обґрунтовано та апробовано концептуальний підхід до цифрової трансформації наукового адміністрування закладу вищої освіти --- переходу від реактивної бібліометричної звітності \emph{post factum} до проактивного управління науковою діяльністю засобами генеративного штучного інтелекту та автономних агентів. Робота позиціонована в полі інформаційно-комунікаційних технологій в управлінні освітою, а не як суто інженерна розробка: технічні механізми тут підпорядковано управлінсько-освітній проблематиці інтеграції розрізнених даних про науковий доробок НПП у єдиний достовірний інформаційний ресурс (research information management).

Отримано такі основні результати. По-перше, розроблено концептуальну модель проактивного управління, у якій ResearchIQ поєднує багатокористувацький вебзастосунок, tenant-сховище даних, серверний API та агентний шар на основі \texttt{pi-agent-core}. Автономний агент інтерпретує природномовні запити керівника, викликає типізовані інструменти аналізу й звітування, але не підміняє детермінований контур прийняття критичних рішень. По-друге, формалізм «золотого запису» не лише описано, а й \emph{реалізовано в коді}: канонізація метаданих, блокування, парне логістичне зіставлення, смуга ручної перевірки, кластеризація із захистом від хибних злиттів за конфліктом DOI та провенанс-злиття атрибутів у дусі W3C~PROV.

По-третє, побудовано відтворюваний синтетичний бенчмарк для реалістичного українського ЗВО (300~НПП, 20~кафедр, приблизно 5~факультетів; без персональних даних), який за процедурою Монте-Карло засвідчив перевагу повного «золотого запису» над базовими методами: $F_1$ дедуплікації зростає до $76{,}8\,\%$ проти $59{,}7\,\%$ у DOI-only та $53{,}5\,\%$ у Title$+$Year, переважно за рахунок повноти ($68{,}1\,\%$ проти $42{,}6\,\%$). Окремо підтверджено критичне значення виправленого дефекту нормалізатора: попередня реалізація знищувала $100\,\%$ україномовних назв. По-четверте, аналітична модель ефективності для 7056 середніх записів оцінює традиційні витрати ручної праці на рівні $\approx 429{,}1$~год/рік на ЗВО ($\approx 21{,}45$~год на кафедральний звіт), тоді як рушій виконує дедуплікацію та злиття за $1{,}58 \pm 0{,}15$~с. По-п'яте, систему спроєктовано за TRACE: компоненти рознесено на детерміноване ядро, класичний ML і мовну модель-валідатор, а довіра до системи є властивістю архітектури, а не окремої моделі \cite{Zabolotnii2026TRACE}.

Теоретичний внесок полягає в концептуалізації наукового адміністрування як самостійного об'єкта цифрової трансформації засобами ІКТ в управлінні освітою та в перенесенні фокуса генеративного ШІ з навчального процесу, де зосереджена переважна частка досліджень, на досі недоопрацьовану адміністративну сферу управління наукою. Практична цінність --- готовий до апробації механізм, що звільняє науково-педагогічних працівників від рутинної верифікації даних, узгоджує звітність із вимогами МОН і створює підґрунтя для євроінтеграційних ініціатив OpenAIRE/EOSC та стандарту CERIF. Відкритий репозиторій містить повний вихідний код ResearchIQ як програмної реалізації запропонованої TRACE-методології (вебзастосунок, серверний агентний шар, детермінований рушій Golden Record), а також скрипти відтворюваного синтетичного бенчмарку, модульні тести й артефакти рукопису; матеріали поширюються під ліцензією MIT: \url{https://github.com/SZabolotnii/ORCID-OpenAlex-Research-Aggregator}.

Перспективи подальших досліджень визначаються й окресленими загрозами валідності. Синтетичні розподіли спотворень є модельними припущеннями про неоднорідність джерел, а не виміряними з них величинами; робастними слід вважати відносні порівняння методів і тренди чутливості, а не абсолютні значення метрик, а оцінку ручного часу --- прозорою аналітичною моделлю, а не вимірюванням за участю людей. Тому першочерговим напрямом є підтвердження результатів на реальній gold-вибірці одного ЗВО з контрольованою участю експертів у контурі (human-in-the-loop). Перспективними є також: розширення агента до рекомендаційних систем добору наукових видань під профіль дослідника; залучення мультимодальних моделей для опрацювання нетекстових артефактів наукової діяльності; розроблення етичних та нормативних рамок застосування автономних агентів в адмініструванні науки, зокрема контролю ризику галюцинацій генеративних моделей у відповідальних управлінських рішеннях.
